
\medskip
\noindent\textbf{4.1.3.}
\phantomsection
\label{sga3:VIIA:4.1.3}

More generally, for every \(S\)-functor \(X\), define the
\(S\)-functor \(X^{(p^n)}\) recursively in \(n\) by
\[
  X^{(p)}=X^{(p/S)}
  \qquad\text{and}\qquad
  X^{(p^n)}
  =
  \bigl(X^{(p^{n-1})}\bigr)^{(p)}.
\]
Likewise, \(\operatorname{Fr}^n(X/S)\), or simply
\(\operatorname{Fr}^n\), denotes the composite functorial
homomorphism
\begin{center}
  \phantomsection
  \label{sga3:VIIA:4.1.3-iterated-relative-frobenius-chain}
  \(\displaystyle
X \xrightarrow{\operatorname{Fr}(X/S)} X^{(p)}
  \xrightarrow{\operatorname{Fr}(X^{(p)}/S)} X^{(p^2)}
  \longrightarrow \cdots \longrightarrow X^{(p^{n-1})}
  \xrightarrow{\operatorname{Fr}(X^{(p^{n-1})}/S)} X^{(p^n)} .
\)

\end{center}
By \SGARefLink{sga3:VIIA:4.1.1}{4.1.1},
\(\operatorname{Fr}(X^{(p)}/S)\) coincides with
\(\operatorname{Fr}(X/S)^{(p)}\); in other words, the following
diagram commutes:
\begin{center}
  \phantomsection
  \label{sga3:VIIA:4.1.3-iterated-frobenius-naturality-square}
  \begin{tikzcd}[row sep=3.2em,column sep=4.2em]
X^{(p)} \arrow[r] \arrow[d,swap,"\operatorname{Fr}(X^{(p)}/S)"]
  & X \arrow[d,"\operatorname{Fr}(X/S)"] \\
X^{(p^2)} \arrow[r]
  & X^{(p)} .
\end{tikzcd}

\end{center}

If \(G\) is a group \(S\)-functor, then \(G^{(p^n)}\) is again a
group \(S\)-functor and
\(\operatorname{Fr}^n(G/S)\) is a homomorphism of group
\(S\)-functors.

\medskip
\noindent\textbf{Definition.}
\phantomsection
\label{sga3:VIIA:4.1.3-height-definition}
\textit{We denote the kernel of \(\operatorname{Fr}^n(G/S)\) by
\(\operatorname{Fr}^nG\), and say that \(G\) has height
\(\leq n\) if \(\operatorname{Fr}^n(G/S)\) is zero, that is, if
\(\operatorname{Fr}^nG=G\).}

\medskip
\noindent\textbf{Lemma.}
\phantomsection
\label{sga3:VIIA:4.1.3-characteristic-kernel-lemma}
\textit{The group subfunctor \(\operatorname{Fr}^nG\) of \(G\) is
characteristic: for every \(S\)-scheme \(T\), every endomorphism
\(\varphi\) of the group \(T\)-functor \(G_T\) induces an
endomorphism of \((\operatorname{Fr}^nG)_T\).}

\begin{proof}
By \SGARefLink{sga3:VIIA:4.1.1}{4.1.1}, the constructions of \(G^{(p^n)}\) and
\(\operatorname{Fr}^n(G/S)\) commute with base change. We may
therefore suppose that \(T=S\). The assertion then follows from the
fact that \(\operatorname{Fr}^n(G/S)\) is a functorial
homomorphism.
\end{proof}

\medskip
\noindent\textbf{4.1.4. Examples.}
\phantomsection
\label{sga3:VIIA:4.1.4}

\smallskip
\noindent\textit{a)}
Let \(M\) first be an ``abstract'' abelian group, and let
\[
  G=D_S(M)
\]
be the diagonalizable group of type \(M\) (I, \SGARefLink{sga3:i:subsection:4-4}{4.4}). For every
\(S\)-scheme \(T\),
\[
  G(T)
  =
  \operatorname{Hom}_{(\mathbf{Ab})}
    \bigl(M,\Gamma(T,\mathcal O_T)^\times\bigr).
\]
Since \(G\) is the pullback of the diagonalizable group \(D(M)\)
over \(\mathbb F_p\), \(G^{(p)}\) identifies with \(G\), and
\(\operatorname{Fr}(G/S)(T)\) identifies with the endomorphism
\(x\mapsto x^p\) of \(G(T)\) (\SGARefLink{sga3:VIIA:4.1.1}{4.1.1}). In particular, when
\(M=\mathbb Z\), one has \(D_S(M)=\mathbb G_{m,S}\). Hence:

\begin{quote}
\(\operatorname{Fr}\mathbb G_{m,S}\) is the \(S\)-group
\(\boldsymbol\mu_{p,S}\) which associates to every \(S\)-scheme
\(T\) the group of \(p\)-th roots of unity in
\(\Gamma(T,\mathcal O_T)^\times\).
\end{quote}

\smallskip
\noindent\textit{b)}
Now let \(S\) be a scheme of characteristic \(p\), and let
\(\mathcal E\) be a sheaf of modules on \(S\). By I, \SGARefLink{sga3:i:proposition:4-6-2}{4.6.2} there is
a canonical isomorphism
\[
  \mathcal W(\mathcal E)^{(p)}
  \xrightarrow{\sim}
  \mathcal W(\mathcal E^{(p)}),
\]
where \(\mathcal E^{(p)}\) is the pullback of \(\mathcal E\) along
\(\operatorname{fr}(S)\). For every \(S\)-scheme
\(\pi:T\to S\), the map
\(\operatorname{Fr}(\mathcal W(\mathcal E)/S)(\pi)\) is determined,
by \SGARefLink{sga3:VIIA:4.1}{4.1}~\((\dagger)\), by the commutative triangle
\begin{center}
  \phantomsection
  \label{sga3:VIIA:4.1.4-vector-group-frobenius-triangle}
  \begin{tikzcd}[row sep=3.6em,column sep=2.6em]
\Gamma\!\left(T,\pi^*\operatorname{fr}(S)^*\mathcal E\right)
  \arrow[rr,"\sim","{\mathrm{can.}}"']
  && \Gamma\!\left(T,\operatorname{fr}(T)^*\pi^*\mathcal E\right) \\
& \Gamma(T,\pi^*\mathcal E)
  \arrow[ul,"\operatorname{Fr}(\mathcal W(\mathcal E)/S)(\pi)"]
  \arrow[ur,"f'"']
& ,
\end{tikzcd}

\end{center}
where \(f'\) is the map induced by \(\operatorname{fr}(T)\).

In particular, if \(\mathcal E=\mathcal O_S\), then
\(\mathcal W(\mathcal E)\) identifies with the additive group
\(\mathbb G_{a,S}\). In this case
\(\mathcal E^{(p)}=\mathcal E=\mathcal O_S\), and the Frobenius
morphism \(\operatorname{Fr}(\mathbb G_{a,S}/S)\) sends
\(x\in\Gamma(T,\mathcal O_T)\) to \(x^p\). Thus:

\begin{quote}
\(\operatorname{Fr}\mathbb G_{a,S}\) is the \(S\)-group
\(\boldsymbol\alpha_{p,S}\) which associates to every
\(S\)-scheme \(T\) the group
\[
  \bigl\{x\in\Gamma(T,\mathcal O_T)\mid x^p=0\bigr\}.
\]
\end{quote}

\smallskip
\noindent\textit{c)}
Similarly, for every quasi-coherent \(\mathcal O_S\)-algebra
\(\mathcal A\), the scheme
\((\operatorname{Spec}\mathcal A)^{(p)}\) identifies with
\(\operatorname{Spec}\mathcal A^{(p)}\), the spectrum of the
pullback of \(\mathcal A\) along \(\operatorname{fr}(S)\). If
\(\pi\) denotes the endomorphism \(x\mapsto x^p\) of the sheaf of
rings \(\mathcal O_S\), then
\[
  \mathcal A^{(p)}
  =
  \mathcal A\otimes_\pi\mathcal O_S,
  \tag{44}\label{sga3:vii:paragraph:4-1-4:equation:44}
\]
and \(\operatorname{Fr}((\operatorname{Spec}\mathcal A)/S)\) is
induced by the morphism of \(\mathcal O_S\)-algebras
\[
  \mathcal A\otimes_\pi\mathcal O_S
  \longrightarrow
  \mathcal A,
  \qquad
  a\otimes_\pi x\longmapsto a^p x.
\]
\footnote{\(\mathcal A\otimes_\pi\mathcal O_S\) denotes the
\(\mathcal O_S\)-algebra obtained by extension of scalars along
\(\pi:\mathcal O_S\to\mathcal O_S\); thus
\(ax\otimes_\pi1=a\otimes_\pi x^p\) and
\(x\cdot(a\otimes_\pi1)=a\otimes_\pi x\).}

Finally, for every quasi-coherent \(\mathcal O_S\)-module
\(\mathcal E\), there are canonical isomorphisms
\[
  \mathcal V(\mathcal E)^{(p)}
  \xrightarrow{\sim}
  \mathcal V(\mathcal E^{(p)})
  \qquad\text{and}\qquad
  \mathcal S(\mathcal E)^{(p)}
  \xrightarrow{\sim}
  \mathcal S(\mathcal E^{(p)}),
\]
where \(\mathcal S(\mathcal E)\) denotes the symmetric algebra of
the \(\mathcal O_S\)-module \(\mathcal E\).

\smallskip
\noindent\textit{d)}
Let \(\mathcal U\) be an \(\mathcal O_S\)-coalgebra (\SGARefLink{sga3:VIIB:3.1:lemma}{3.1}), and let
\(S\) be a scheme of characteristic \(p\). If
\(\mathcal U^{(p/S)}\), or \(\mathcal U^{(p)}\), denotes the
pullback of the coalgebra \(\mathcal U\) along
\(\operatorname{fr}(S)\), then, as in \textit{b)}, there is a
canonical isomorphism
\[
  \bigl(\operatorname{Spec}^{*}\mathcal U\bigr)^{(p)}
  \xrightarrow{\sim}
  \operatorname{Spec}^{*}\mathcal U^{(p)}.
\]
If \(\mathcal U\) is a group coalgebra, then for an \(S\)-scheme
\(T\), the value of
\(\operatorname{Fr}(\operatorname{Spec}^{*}\mathcal U)\), that is,
of the kernel of the Frobenius morphism
\[
  \operatorname{Spec}^{*}\mathcal U
  \longrightarrow
  \bigl(\operatorname{Spec}^{*}\mathcal U\bigr)^{(p)},
\]
is the set of elements \(\gamma\) of
\[
  \bigl(\operatorname{Spec}^{*}\mathcal U\bigr)(T)
  =
  \left\{
    x\in\Gamma(T,\mathcal U_T)
    \,\middle|\,
    \varepsilon_{\mathcal U_T}(x)=1,\ 
    \Delta_{\mathcal U_T}x=x\otimes x
  \right\}
\]
such that the image in
\(\Gamma(T,\mathcal U_T\otimes_{\operatorname{fr}(T)}\mathcal O_T)\)
of the element
\[
  \gamma\otimes_{\operatorname{fr}(T)}1
  \in
  \Gamma(T,\mathcal U_T)\otimes_{\operatorname{fr}(T)}
    \mathcal O(T)
\]
is equal to \(1\).

